% Conclusion -- handbook section 2.3 requires this coverage
% What the comparison establishes, the limits, and what comes next.

\section{Conclusion}

%----------------------------------------------------------
\begin{frame}{Findings}
\small
\begin{block}{Research question: does inductive bias or capacity govern performance here?}
\textbf{The evidence favours inductive bias.}
\end{block}

\begin{itemize}
    \item The most accurate model is the \emph{smallest} of the three strong ones, and keeps its margin
          even over the best sequence configuration the ablations found.
    \item The largest model, trained at twice the cost per epoch, is indistinguishable from a cheaper one.
    \item The error structure that dominates every model's residual was predicted from the class means
          before any network was trained.
    \item \textbf{Selected model: the CNN}, $0.9110 \pm 0.0007$ macro-F1 on test, and also the cheapest
          to run among the models competitive on accuracy, so there is no accuracy-against-speed
          trade-off to negotiate here.
    \item \textbf{A result can hinge on a choice made before any measurement:} the Transformer's
          apparent tie with the MLP disappeared once the representation was varied.
\end{itemize}
\end{frame}

%----------------------------------------------------------
\begin{frame}{Limitations and Future Work}
\small
\begin{columns}[T]
\begin{column}{0.5\textwidth}
\begin{block}{Limits}
\begin{itemize}
    \item One dataset of $28\times28$ grayscale images with balanced classes; neither the ranking nor
          the error structure should be assumed to transfer.
    \item Three seeds estimate a deviation from three observations, so a tie means \emph{not separated
          at this sample size}, not equal.
    \item One shared protocol and an equal search cap, so each figure is a score under a common
          budget, not the best each architecture could reach.
    \item No pretrained weights, which constrains the Transformer most.
\end{itemize}
\end{block}
\end{column}
\begin{column}{0.5\textwidth}
\begin{block}{Next}
\begin{itemize}
    \item Ablate the representation for the remaining three architectures; only two of the five were
          tested at more than one view of their input.
    \item Reweight the loss towards Shirt, the single class that bounds macro-F1 for all five models.
    \item Equalise parameter counts across families to separate capacity from architecture directly.
\end{itemize}
\end{block}
\end{column}
\end{columns}
\end{frame}
